\subsection{ADDITIVE UNIFORMITY OF R}

Since the group \(\mathbf{R}\) is linearly ordered, the functions \(x^{+}, x^{-}\) and \(|x|\) are defined on \(\mathbf{R}\) in the same way as on \(\mathbf{Q}\) and satisfy all the relations listed above for \(\mathbf{Q}\) , notably relations (1) to (7). Let \(a\) be a real number \(>0\) and let \(\mathrm{U}_{a}\) be the set of all pairs \((x, y) \in \mathbf{R} \times \mathbf{R}\) such that \(|x - y| < a\) ; as \(a\) runs through the set of real numbers \(>0\) (or the set of numbers \(1/n\) ), the sets \(\mathrm{U}_{a}\) form a fundamental system of entourages of the uniformity of the additive group \(\mathbf{R}\) of the real line (called the additive uniformity of the real line).

The functions \(|x|\) , \(x^{+}\) and \(x^{-}\) are uniformly continuous on R, and the functions sup \((x, y)\) and inf \((x, y)\) are uniformly continuous on \(R \times R\) ; these functions therefore coincide with those obtained by extending by continuity the corresponding functions defined on Q and \(Q \times Q\) , respectively.

